\subsection{Homogene lineare Gleichungssysteme}
\label{subsec:linear_systems_homogeneous_systems}

Ein wichtiger Spezialfall entsteht, wenn sämtliche Absolutglieder eines LGS
null sind.

\begin{DefinitionBox}{Homogenes lineares Gleichungssystem}{linear_systems_homogeneous_system}
  Ein lineares Gleichungssystem heißt \emph{homogen}
  \EnglishTerm{homogeneous}, wenn alle Absolutglieder gleich $\MathHighlight{\text{null}}$ sind:
  \[
    \begin{aligned}
      a_{11}x_1+a_{12}x_2+\dots+a_{1n}x_n&=\MathHighlight{0},\\
      a_{21}x_1+a_{22}x_2+\dots+a_{2n}x_n&=\MathHighlight{0},\\
      &\mathrel{\phantom{=}}\vdots\\
      a_{m1}x_1+a_{m2}x_2+\dots+a_{mn}x_n&=\MathHighlight{0}.
    \end{aligned}
  \]
  Seine erweiterte Koeffizientenmatrix besitzt die Form $[A\mid 0]$.
\end{DefinitionBox}

\begin{TheoremBox}{Triviale Lösung homogener LGS}{linear_systems_homogeneous_trivial_solution}
  Jedes homogene LGS besitzt die \emph{triviale Lösung}
  \EnglishTerm{trivial solution}
  \[
    (x_1,\dots,x_n)=\Cloze{4.5cm}{(0,\dots,0)}.
  \]
  Denn beim Einsetzen dieser Werte wird jede Gleichung zu $0=0$.
\end{TheoremBox}

\begin{KeyPointBox}{Lösungsmengen homogener LGS}{linear_systems_homogeneous_solution_sets}
  Da die triviale Lösung immer vorhanden ist, kann die Lösungsmenge eines
  homogenen LGS nicht leer sein:
  \SolutionTypeCases*
  Ein homogenes LGS besitzt daher entweder genau eine oder unendlich viele
  Lösungen.
\end{KeyPointBox}
